\input{_preamble}
%% \AltJdgTuck is not clamped, and a tuck too deep is a SILENT fault: TeX
%% has no opinion about overlapping ink, so the mark simply prints on top
%% of the brace with no error and nothing in the log.  The package checks
%% the value and reports rather than correcting it -- a tunable that
%% quietly ignores what it is set to would be worse than one that obeys.
%%
%% Asserted on the .log, since the warning is the whole point; the geometry
%% of a well-behaved tuck is pinned in altn-phantomalign.tex instead.
\GlossPhantomAlign
\begin{document}

% The default must be silent.  It sits below the threshold and is the
% value every document that never touches the knob will use, so a warning
% here would be a warning for everybody.
\noindent QUIETL \lxAltn[l]{QSTEM}{*QSTEM} QUIETR

% A stack with no mark does not tuck at all, so however absurd the length
% is it cannot collide -- and must not be reported.
{\renewcommand\AltJdgTuck{40pt}
\noindent NOMARKL \lxAltn[l]{NSTEM}{NSTEM} NOMARKR\par}

% Past \AltBraceWidth + \AltBraceSep the stack is pulled beyond the brace's
% own box.  This one is reported.
{\renewcommand\AltJdgTuck{18pt}
\noindent DEEPL \lxAltn[l]{DSTEM}{*DSTEM} DEEPR\par}

% ... and reported ONCE, not once per stack: an inline construction in a
% long paper would otherwise fill the log with copies of itself.
{\renewcommand\AltJdgTuck{22pt}
\noindent AGAINL \lxAltn[l]{ASTEM}{*ASTEM} AGAINR\par}
\end{document}
